\part{Markets, Firms, Commons, and Coordination}
\label{part:VII}

Part VII asks when individually acceptable commitments compose into a feasible collective arrangement, in markets, firms, platforms, and commons.

\begin{housing}A shared crew serves several buildings, listings run through a platform, and a cooperative maintenance fund pools contributions.\end{housing}

\chapter{Exchange and Compatible Commitments}\label{ch:31}
Models exchange as a transformation of resources, permissions, and obligations among participants. Identifies conditions under which individually acceptable exchanges compose into a feasible collective arrangement.
\chspec{\Pref{P09}, \Pref{P10}, \Pref{P18}.}
\draftnote{Candidate case (documentation-level only): two customers each receive an apparently feasible offer for the last item while one purchase can be fulfilled. Reservations add time: resources are withheld, expire, and return. Authorization and resource commitment must stay consistent under concurrent requests (P09).}
\scaffold

\chapter{Firms as Organizations of Complementary Continuation}\label{ch:32}
Treats firms as coordinating tasks, assets, knowledge, and promises that cannot be maintained through isolated transactions alone. Boundaries of the firm are examined through coordination costs, authority, dependence, and the distribution of residual responsibilities.
\chspec{\Pref{P27}.}
\scaffold

\chapter{Platforms and Infrastructural Dependence}\label{ch:33}
Studies platforms as institutions that establish opportunity domains while controlling access to them. Examines interoperability, portability, unilateral rule changes, and the appropriation of complementary investment.
\draftnote{Candidate import: revenue with a participation term, $R(q)=(\alpha q+\beta)M(q)$, so $R'(q)=\alpha M(q)+(\alpha q+\beta)M'(q)$. Revenue can fall as match quality rises if participation falls strongly enough with quality. This is a conditional incentive result only.}
\scaffold

\chapter{Commons and Shared Maintenance}\label{ch:34}
Analyzes shared resources through contribution rules, monitoring, access, replenishment, and conflict resolution. Distinguishes common ownership from the actual organization of maintenance.
\scaffold

\chapter{Contribution, Reward, and Measurement}\label{ch:35}
Uses work-backed currencies, contribution scores, and automated reward systems as worked cases. Shows how reward mechanisms change the activity they measure and why maintenance and preventive work are particularly hard to recognize through simple activity counts.
\chspec{\Pref{P15}.}
\scaffold
